\section{Preliminary results}
\label{sec_pre}
Throughout this paper, we refer to \cite[Chapter 2]{KM98} for different kinds of singularities.
By a \textit{projective klt (resp. canonical, dlt) pair} $(X,\Delta)$, we mean that $X$ is a normal projective variety, $\Delta\geqslant 0$ is an effective $\mathbb{Q}$-divisor such that $K_X+\Delta$ is $\mathbb{Q}$-Cartier, and  $(X,\Delta)$ has only klt (resp. canonical, dlt) singularities.

A $\mathbb{Q}$-Cartier divisor $L$ on a projective variety is said to be \textit{strictly nef} if $L\cdot C>0$ for every (complete) curve $C$ on $X$.
We also recall the definition of almost strictly nef divisors for the convenience of the proofs in the later sections.
\begin{defn}[{cf.~\cite[Definition 1.1]{CCP08}}]\label{defn-almost-sn}
A $\mathbb{Q}$-Cartier divisor $L$ on a normal projective variety $X$ is called  \textit{almost strictly nef}, if there exist a birational morphism $\pi:X\to Y$ to a  normal projective variety $Y$ and a strictly nef $\mathbb{Q}$-divisor $L_Y$ on $Y$ such that $L=\pi^*L_Y$.  
\end{defn}
One of the motivation to introduce almost strictly nef divisors is to study the  non-uniruled case of {\hyperref[main-conj-singular-arbitrary]{Question \ref*{main-conj-singular-arbitrary}}}  via the descending of Iitaka fibrations, in which case, the pullback of a strictly nef divisor to a higher model which resolves the indeterminacy is almost strictly nef  (cf.~\cite[Theorem 2.6]{CCP08}). 


Now, we show  that, under the condition of  {\hyperref[main-conj-singular]{Question \ref*{main-conj-singular}}},  such $X$ is uniruled.
This permits us to take a first glance on the geometric property of a  projective klt pair with (almost) strictly nef anti-log canonical divisor. 
%We also note that, {\hyperref[main-conj-singular]{Question \ref*{main-conj-singular}}} has an affirmative answer when $X$ is smooth with the boundary $\Delta=0$ (cf.~ \cite[Theorem 1.3]{LOY19}).


\begin{prop}\label{prop_strict_uniruled}
Let $(X,\Delta)$ be a projective pair such that $-(K_X+\Delta)$ is almost strictly nef.
Then $X$ is uniruled, i.e., $X$ is covered by rational curves.	
\end{prop}
\begin{proof}
Suppose the contrary that $X$ is not uniruled.
Then $K_X$ is pseudo-effective as a Weil divisor (cf. \cite{BDPP13}) by considering a resolution of $X$.
In other words, $K_X$ lies in the closure of the cone generated by effective Weil divisors on $X$ (cf.~\cite[Definition 2.2]{MZ18}).   
%\cite[Chapter II, Definition 5.5]{Nak04} and
So $-\Delta=-(K_X+\Delta)+K_X$ is also pseudo-effective (as a 
Weil divisor) and thus $\Delta=0$.
This in turn implies the almost strict nefness of $-K_X$, a contradiction to the pseudo-effectivity of $K_X$. 
\end{proof}




In the remaining part of this section, we review the definition and basic properties of locally constant fibrations and flat vector bundles for the convenience of readers.

Let $\phi:X\rightarrow Y$ be a proper morphism between normal varieties. We say that $\phi$ is a \emph{fibration} if $\phi_*\scrO_X=\scrO_Y$. 
\begin{defn}[{\cite[Definition 2.6]{MW21}; cf.~\cite[Definition 1.6]{Wang20}}]\label{defn-locally-constant}
Let $\phi:X\to Y$ be a surjective morphism with connected fibres between analytic varieties, and let $\Delta$ be a Weil $\mathbb{Q}$-divisor on $X$.
\begin{enumerate}
	\item[(1)] The morphism $\phi:X\to Y$ is said to be a \textit{locally constant fibration with respect to the pair} $(X,\Delta)$ if it satisfies the following conditions:
	\begin{itemize}
	\item The morphism $\phi:X\to Y$  is a (locally trivial) analytic fibre bundle with the fibre $F$.
	\item Every component $\Delta_i$ of $\Delta$ is horizontal (i.e., $\phi(\Delta_i)=Y$).
	\item There exist a representation
	$$\rho:\pi_1(Y)\to \textup{Aut}(F)$$
	of the fundamental group $\pi_1(Y)$ of $Y$ to the automorphism group $\Aut(F)$ of $F$ and a Weil $\mathbb{Q}$-divisor $\Delta_F$ on $F$ which is invariant under the action of $\pi_1(Y)$, so that $(X,\Delta)$ is isomorphic to the quotient of $(Y^{\text{univ}}\times F, \textup{pr}_2^*\Delta_F)$ over $Y$ by the action of $\pi_1(Y)$ on $Y^{\textup{univ}}\times F$.
	Here $Y^{\textup{univ}}$ is the universal cover of $Y$, and the action of $\gamma\in\pi_1(Y)$ is defined to be
	$$\gamma\cdot(y,z):=(\gamma\cdot y,\rho(\gamma)(z))$$
	for any $(y,z)\in Y^{\textup{univ}}\times F$.
	\end{itemize}
 \item[(2)] The morphism $\phi:X\to Y$ is simply said to be \textit{locally constant} if it satisfies the above conditions for $\Delta=0$.
 \item[(3)] For the convenience of the notation, we denote by $\text{Aut}(F,\Delta_F)$  the group of  automorphisms of $F$ which leave  $\Delta_F$ invariant, counted with the multiplicities.
\end{enumerate}
\end{defn}




One of the most important examples of locally constant families  are   \emph{flat vector bundles}, which plays an important role in this paper.

\begin{defn}
Let $E\rightarrow Y$ be a holomorphic vector bundle of rank $r$ over a normal variety. Then we say that $E$ is \textit{flat} if $E\rightarrow Y$ is a locally constant family corresponding to a representation $\rho:\pi_1(Y)\rightarrow \text{GL}(r,\mathbb{C})$.
\end{defn}

Let $E\rightarrow Y$ be a holomorphic vector bundle over a normal projective variety.
Recall that $E$ is called \emph{nef} if the tautological line bundle $\scrO_{\PP E}(1)$ over $\PP E$ is nef, and $E$ is called \emph{numerically flat} if both $E$ and its dual bundle $E^*$ are nef. According to \cite[Theorem 1.18 and Corollay 1.19]{DPS94} and \cite[Section 3]{Sim92},  a numerically flat vector bundle $E\rightarrow Y$ over a normal projective variety is a flat vector bundle. 


\begin{lemme}[{\cite[Lemma 4.4]{LOY19}}]
\label{lemma_fixed-pt-subbundle}
Let $V$ be a finite dimensional representation of $G:=\pi_1(Y)$ with $E$ being the flat vector bundle on $Y$ induced by this representation. 
Then there is a one-to-one correspondence between the set of $G$-fixed points $z\in\PP V$ and the set of codimension one flat subbundles of $E$.
\end{lemme}


This is nothing but a tautology fact (functorial definition of projective spaces); see \cite[Lemma 4.4]{LOY19} for the case when $Y$ is a smooth projective variety. 
Indeed, the projectivity and smoothness of  $Y$ are not needed in the proof. 
For readers' convenience, we include a detailed proof.

\begin{proof}[Proof of {\hyperref[lemma_fixed-pt-subbundle]{Lemma \ref*{lemma_fixed-pt-subbundle}}}]
Clearly, there is a one-to-one correspondence between the set of $G$-fixed-points $z\in\PP V$ (which corresponds to a hyperplane $V_z\subset V$) and the set of quotient representations $V\to V/V_z$.
By the bijection between the isomorphism classes of  local systems on $Y$ and isomorphism classes of  representations of the fundamental group $\pi_1(Y)$, there is a one-to-one correspondence between  the set of $G$-fixed-points $z\in\PP V$ and line bundle quotients $E\to F$.
Since the kernel of a surjective morphism between locally free sheaves is locally free, our lemma follows.
\end{proof}






\begin{lemme}[{\cite[Lemma 4.5]{LOY19}}]
	\label{l.fixedpointflatsection}
	Let $Y$ be a normal projective variety and $\pi:Y\univ\rightarrow Y$  the universal cover with Galois group $G=\pi_1(Y)$. 
	Let $\rho:G\rightarrow \text{GL}(V,\CC)$ be a linear representation of $G$ and denote by $E$ the corresponding flat vector bundle over $Y$. 
	If $z\in\PP(V)$ is a $G$-fixed-point, then it induces a section $\sigma:Y\rightarrow \PP E$   such that $\sigma^*\scrO_{\PP E}(1)$ is numerically trivial. 
\end{lemme}

\begin{proof}
Let $F$ be  the co-rank one flat  subbundle of 
$E$ corresponding to $z$ (cf.~{\hyperref[lemma_fixed-pt-subbundle]{Lemma \ref*{lemma_fixed-pt-subbundle}}}).  
Let $\sigma:Y\rightarrow \PP E$ be the section defined by the following short exact sequence
	\[
	0\rightarrow F\rightarrow E\rightarrow E/F\rightarrow 0.
	\]
	Then we get $\sigma^*\scrO_{\PP E}(1)\cong E/F$. 
	On the other hand, as both $F$ and $E$ are numerically flat, it follows that $E/F$ and hence $\sigma^*\scrO_{\PP E}(1)$ are also numerically flat.
In particular, the line bundle $\sigma^*\scrO_{\PP E}(1)$ is numerically trivial.
\end{proof}

\begin{rmq}
	In the above lemma, the image $\sigma(Y)$ is just the image of $Y\univ\times \{z\}$ under the natural morphism $Y\univ\times \PP(V)\rightarrow \PP E$. 
	Here, we identify $\PP E$ with  the quotient of $Y\univ\times \PP V$ by the induced diagonal action of $G$.
\end{rmq}



Let $f:X\rightarrow Y$ be a locally constant fibration with respect to the pair $(X,\Delta)$, and let $\mu:Y'\rightarrow Y$ be a morphism from a normal variety $Y'$. 
Then there exists a locally constant fibration $f':X'\rightarrow Y'$ with respect to a pair $(X',\Delta')$ induced by $f$. Indeed, we consider the natural composition
\[
\pi_1(Y')\xrightarrow{\mu_*} \pi_1(Y) \xrightarrow{\rho} \Aut(F,\Delta_F), 
\]
where $\rho$ is the representation associated to $f$. Then   $X' = X\times_{Y} Y'$.

With this base change kept in mind, we prove the following lemma.
\begin{lemme}
	\label{l.basechange}
	Let $f:X\rightarrow Y$ be a locally constant fibration between normal projective varieties with respect to a pair $(X,\Delta)$. Assume that $K_X+\Delta$ is $\QQ$-Cartier and $Y$ is $\QQ$-Gorenstein. Let $\mu:Y'\rightarrow Y$ be a morphism from a normal projective variety $Y'$ with $\QQ$-Gorenstein singularities. 
	Denote by $f':X'\rightarrow Y'$ the induced locally constant fibration by $f$ and let $g:X'\rightarrow X$ be the natural   morphism. 
	Then we have
	\[
	g^*(K_{X/Y}+\Delta) \sim_{\QQ} K_{X'/Y'}+\Delta'.
	\]
\end{lemme}

\begin{proof}
	Let $F$ be a general fibre of $f$ and let $\Delta_F$ be the $\QQ$-Weil divisor on $F$ as  in {\hyperref[defn-locally-constant]{Definition  \ref*{defn-locally-constant}}}. 
	Let $Y\univ$ and $(Y')\univ$ be the universal coverings of $Y$ and $Y'$, respectively. Then the morphism $\mu:Y'\rightarrow Y$ can be lifted to a morphism $\widetilde{\mu}:(Y')\univ\rightarrow Y\univ$ satisfying the following commutative diagram
	\[
	\begin{tikzcd}[column sep=large, row sep=large]
		(Y')\univ\times F \arrow[r,"\widetilde{g}"] \arrow[d,"\pi'" left]  
		    & Y\univ \times F  \arrow[d,"\pi"]  \\
		X' \arrow[r,"g" below]
		    & X
	\end{tikzcd}
	\]
	where $\widetilde{g}:=(\widetilde{\mu},\textup{id})$. Since $K_X+\Delta$ is $\QQ$-Cartier, the $\QQ$-Weil divisor $K_F+\Delta_F$ is also $\QQ$-Cartier. Moreover, since $\pi'$ and $\pi$ are both \'etale, we get
	\begin{center}
		$(\pi')^*(K_{X'/Y'}+\Delta')=(p_2')^*(K_F+\Delta_F)$ \quad and \quad $\pi^*(K_{X/Y}+\Delta)=p_2^*(K_F+\Delta_F)$,
	\end{center}
    where $p_2'$ and $p_2$ are the projections $(Y')\univ\times F\rightarrow F$ and $Y\univ\times F\rightarrow F$, respectively. 
    In particular, we obtain
    \[
    (g\circ \pi')^*(K_{X/Y}+\Delta) \sim (\pi\circ \widetilde{g})^*(K_{X/Y}+\Delta) \sim (\pi')^*(K_{X'/Y'}+\Delta').
    \]
    Note that, by our assumption, the line bundle $(p'_2)^*\scrO_{F}(mK_F+m\Delta_F)$ is equivariant with respect to the induced action of $\pi_1(Y')$ on $(Y')\univ\times F$ for 
     for $m$ sufficiently divisible. Thus we have $g^*(K_{X/Y}+\Delta)\sim K_{X'/Y'}+\Delta'$. 
\end{proof}


\vskip 2\baselineskip

















